% Taxer le terrain : note unique sur la taxe sur la valeur foncière (land value tax)
% Compilation : pdflatex lvt-brief.tex (deux passes) depuis fiscalite/lvt/
\documentclass[11pt,a4paper]{article}
% Préambule commun au brief et à la revue de littérature sur la taxe sur la valeur foncière.
% Même composition que le cours (cours/refaire-l-impot.tex) ; figures « nues » produites par
%   FIG_NU=1 python lvt/graphiques_lvt.py
\usepackage[T1]{fontenc}
\usepackage[utf8]{inputenc}
\usepackage{lmodern}
\usepackage[french]{babel}
\usepackage[a4paper,left=2.8cm,right=2.8cm,top=2.6cm,bottom=2.8cm]{geometry}
\usepackage{graphicx}
\usepackage[export]{adjustbox}
\usepackage{booktabs}
\usepackage{tabularx}
\usepackage{array}
\usepackage{xcolor}
\usepackage{titlesec}
\usepackage[labelfont=bf,labelsep=period,font=small,singlelinecheck=false,
            justification=raggedright,skip=6pt]{caption}
\usepackage{fancyhdr}
\usepackage{enumitem}
\usepackage{microtype}
\usepackage{textcomp}
\usepackage{amsmath}
\providecommand{\euro}{\texteuro}
\usepackage[hidelinks]{hyperref}

\definecolor{encre}{HTML}{1A1A19}
\definecolor{bleu}{HTML}{003399}
\definecolor{vert}{HTML}{1F7A4D}
\definecolor{ocre}{HTML}{C79100}
\definecolor{brique}{HTML}{993333}
\definecolor{filet}{HTML}{C8C6BE}

\titleformat{\section}{\normalfont\large\bfseries\color{bleu}}{\thesection.}{0.6em}{}
\titleformat{\subsection}{\normalfont\normalsize\bfseries\color{bleu}}{\thesubsection}{0.6em}{}
\titlespacing*{\section}{0pt}{2.4ex plus 1ex minus .2ex}{1.2ex plus .2ex}
\titlespacing*{\subsection}{0pt}{2ex plus 1ex minus .2ex}{0.8ex plus .2ex}

\pagestyle{fancy}
\fancyhf{}
\renewcommand{\headrulewidth}{0.4pt}
\fancyhead[R]{\small\thepage}

\setlength{\parindent}{0pt}
\setlength{\parskip}{0.62em}
\linespread{1.05}
\renewcommand{\arraystretch}{1.2}
\setlist{itemsep=2pt,topsep=2pt,parsep=0pt,leftmargin=1.4em}

\newenvironment{encadre}[1]
  {\par\medskip\noindent\begin{minipage}{\linewidth}%
   \color{encre}\setlength{\parskip}{0.5em}%
   \hrule height 0.8pt width \linewidth \vspace{6pt}%
   {\footnotesize\bfseries\color{vert}\MakeUppercase{#1}}\par\vspace{-2pt}%
   \small}
  {\par\vspace{4pt}\hrule height 0.4pt width \linewidth\end{minipage}\par\medskip}

\newcommand{\figbloc}[3]{%
  \begin{figure}[htbp]\centering
  \caption{#2}\label{#3}
  \includegraphics[max width=\linewidth,max height=0.40\textheight]{../figures/nu/#1.pdf}
  \IfFileExists{../figures/nu/#1_note.tex}{\par\vspace{4pt}%
    \parbox{\linewidth}{\raggedright\scriptsize\color{encre}\setlength{\parskip}{2pt}%
    \input{../figures/nu/#1_note.tex}}}{}%
  \end{figure}}

\newcommand{\pc}{\,\%}

\fancyhead[L]{\small\itshape Taxer le terrain : la taxe sur la valeur foncière}

\begin{document}
\renewcommand{\figurename}{Figure}\renewcommand{\tablename}{Tableau}

\thispagestyle{empty}
\vspace*{0.6cm}
{\footnotesize\color{ocre}\bfseries NOTE DE TRAVAIL\quad·\quad OCTOBRE 2026\par}
\vspace{8pt}
{\fontsize{24}{28}\selectfont\bfseries\color{bleu}Taxer le terrain\par}
\vspace{6pt}
{\large Ce qu'il faut savoir sur la taxe sur la valeur foncière (\emph{land value tax}) et ce
qu'elle donnerait en France\par}
\vspace{8pt}
{\small Colin Baget. Les chiffres français sont calculés sur les comptes de patrimoine de l'Insee,
les transactions DVF 2024 et les données de l'Insee Analyses n°~91 ; la méthode est en annexe.\par}
\vspace{6pt}
{\color{bleu}\hrule height 1.2pt}
\vspace{4pt}

\section*{L'essentiel}

Une taxe sur la valeur foncière (LVT) est un impôt annuel sur la valeur du seul terrain, estimée
comme s'il était nu, quel que soit ce qu'on a construit dessus. Elle séduit les économistes pour
une raison simple : la quantité de terrain ne dépend pas de l'impôt. En théorie, l'impôt ne déforme donc aucune décision et pèse sur les
propriétaires au moment de son annonce, par une baisse du prix du sol.

Les terrains sous les bâtiments sont passés de 53~\% du PIB en 1998 à 229~\% fin 2024. Ce chiffre
est un résidu statistique fragile, mais la tendance est sûre : les prix des logements ont triplé depuis 1998 quand le coût de la construction doublait et
les loyers suivaient l'inflation. La France a aussi une taxe foncière (42~Md€) assise sur des
valeurs locatives de 1970, qui pèse 0,54~\% de la valeur des logements dans les communes les moins
chères contre 0,20~\% dans les plus chères.

Les preuves empiriques sont plus minces que la théorie. Transposées à la France, les études
pennsylvaniennes donnent de 8~000 à 60~000 logements de plus par an, un appoint plutôt qu'un
choc d'offre. Sur la question décisive, savoir qui paie, les deux études de la même réforme
danoise se contredisent : l'une trouve que la taxe est entièrement capitalisée dans les prix,
l'autre, plus récente, qu'elle ne l'est pas du tout, auquel cas elle pèse sur les acheteurs
successifs et les locataires. Les preuves les plus solides portent sur les droits de mutation,
que la LVT pourrait remplacer : ils freinent nettement ventes et déménagements à l'étranger, alors
qu'en France la hausse de 2014 n'a eu qu'un effet de calendrier.

Une taxe sur le terrain remplaçant la taxe foncière au même rendement demanderait environ 0,7~\% de
la valeur des terrains par an, 1~\% si elle remplaçait aussi les droits de mutation. Elle
déplacerait fortement la charge vers les zones chères : une maison médiane de Corrèze et une de
Haute-Savoie paient aujourd'hui la même taxe ; la seconde paierait 25 fois plus que la première.
Étaler la réforme réduit peu les pertes en capital si les acheteurs anticipent : avec une montée
en charge sur dix ans, 88~\% de la perte tombe le jour de l'annonce. Comme ils anticipent mal, en
pratique, l'étalement reporte une partie de la charge sur ceux qui achètent pendant la
transition.

Pour la France, la question est de changer l'assiette d'un impôt existant. Une trajectoire
réaliste commence par asseoir la taxe foncière sur des valeurs de marché, puis accroît
la part du terrain dans l'assiette, en finançant une baisse des droits de mutation. Le vrai
obstacle est l'estimation de la valeur de chaque terrain, que personne ne fait aujourd'hui en
France.

\section{De quoi parle-t-on}

Un bien immobilier réunit deux choses de nature différente : un bâtiment, que quelqu'un a produit,
qui s'use et qu'on peut reconstruire, et un emplacement, que personne n'a produit et qui ne s'use
pas. La LVT ne taxe que le second, sur la valeur qu'aurait le terrain s'il était vide et disponible
pour l'usage le plus rentable que l'urbanisme autorise. Un parking à ciel ouvert et un immeuble de
six étages sur deux parcelles voisines de même surface paient donc le même impôt. Un impôt foncier
ordinaire fait l'inverse : en France, le parking paie la taxe sur le non-bâti ou une taxe bâtie
faible, l'immeuble paie la taxe foncière sur chacun de ses logements ; construire alourdit
l'impôt, laisser en friche l'allège. Entre les deux, la taxe à deux taux (\emph{split-rate})
applique au terrain un taux plus élevé qu'au bâti, comme Pittsburgh de 1913 à 2001 et une vingtaine
de villes de Pennsylvanie. La France ne taxe le sol en tant que tel qu'au moment où il devient
constructible (taxe communale facultative de 10~\% de l'article 1529 du CGI, taxe nationale de
l'article 1605~\emph{nonies}).

Selon l'OCDE (2022), seuls le Danemark (taxe communale de 16 à 34~‰ de la valeur du sol de 2007 à
2023), l'Estonie et la Nouvelle-Galles du Sud (qui exonère la résidence principale) taxent le
terrain seul.

L'expérience la plus proche est allemande. En 2018, la Cour constitutionnelle a censuré l'impôt
foncier, assis sur des valeurs de 1964 à l'Ouest et de 1935 à l'Est, pour le défaut qui rend la
taxe foncière française inéquitable : des valeurs figées traitent différemment des biens
semblables. La réforme appliquée en 2025 a laissé chaque Land choisir. Le Bade-Wurtemberg taxe le
seul terrain : surface de la parcelle multipliée par la valeur de référence du sol de la zone
(\emph{Bodenrichtwert}), publiée par des comités d'experts à partir des ventes ; taux de base de
1,30~‰ réduit de 30~\% pour le résidentiel ; preuve d'une valeur inférieure admise au-delà de
30~\% d'écart. La Cour fédérale des finances a jugé le 20~mai 2026 que ce modèle n'était pas
contraire à la Constitution. Les communes peuvent aussi taxer plus lourdement les terrains
constructibles laissés vides (\emph{Grundsteuer~C}).

\section{Ce que dit la théorie}

\subsection{Un impôt que personne ne peut éviter}

Le coût d'un impôt vient des décisions qu'il fait changer. Le terrain est le cas limite d'une assiette qui ne réagit pas : taxé, il ne peut ni être
produit en moindre quantité, ni partir, ni être caché. Son propriétaire ne peut pas non plus
reporter l'impôt sur ses locataires, car le loyer d'un emplacement est fixé par ce que les
candidats sont prêts à payer pour être là plutôt qu'ailleurs ; l'impôt ne change ni le nombre
d'emplacements, ni leur attrait, ni le nombre de candidats.

Ce que perd le propriétaire se lit dans le prix. Un terrain qui rapporte un loyer net de 3~000~€
par an vaut 100~000~€ si l'on actualise à 3~\%. Une taxe annuelle de 1~\% de sa valeur est annoncée :
l'acheteur ne paiera plus que le prix $P$ tel que $3\,000 - 0{,}01\,P = 0{,}03\,P$, soit 75~000~€.
Le propriétaire perd 25~000~€ le jour de l'annonce, la valeur actualisée de tous les impôts futurs ;
l'acheteur suivant paie 750~€ d'impôt par an mais a payé 25~000~€ de moins. C'est la
capitalisation : la taxe fonctionne comme un prélèvement unique sur le patrimoine foncier existant,
que rien ne permet d'éviter, ce qui la rend efficace et politiquement coûteuse
(figure~\ref{fig:capitalisation}). La \emph{Mirrlees Review} juge l'argument « simple et presque
indéniable ».

\figbloc{lvt_capitalisation}{Effet d'une taxe annuelle sur le prix du terrain, selon le taux
d'actualisation}{fig:capitalisation}

Deux résultats théoriques nuancent ce schéma. Si les ménages épargnent pour leur retraite en
détenant terre ou capital, la taxe réoriente l'épargne vers le capital, fait baisser le taux
d'intérêt et monter les salaires : le prix du sol baisse alors moins que prévu (Feldstein, 1977),
sauf si les ménages lèguent par altruisme (Calvo, Kotlikoff et Rodriguez, 1979). Et la neutralité dépend de
l'évaluation. Taxé à son usage courant, un terrain qui ne rapportera qu'après aménagement paie dès
aujourd'hui sur des revenus futurs, ce qui pousse vers les usages qui rapportent tout de suite
(Bentick, 1979 ; Mills, 1981) ; évalué à son meilleur usage permis, quel que soit ce qu'en fait le
propriétaire, il ne déforme plus rien (Tideman, 1982 ; Wildasin, 1982).

\subsection{Ce que la taxe ne fait pas : obliger à construire}

On lit souvent que la LVT « force » à mettre en valeur les terrains vides. C'est inexact : comme
l'impôt est le même quel que soit l'usage, il ne change le classement d'aucun projet (Oates et
Schwab, 1997). L'effet favorable à la construction vient de ce qu'on retire en même temps : un
impôt sur le bâti qui, lui, décourage de construire. Une taxe qui vise la non-utilisation
elle-même produit en revanche un effet : la taxe sur les logements vacants de 1999 a réduit la
vacance de 13~\% dans les agglomérations concernées (Segú, 2020). La densification n'est pas acquise non plus :
l'impôt sur la propriété réduit le capital investi par hectare, ce qui étale la ville, mais aussi la
surface de logement consommée, ce qui la densifie (Brueckner et Kim, 2003) ; le supprimer au profit
d'une LVT a donc un effet ambigu, que seules les données tranchent.

\subsection{La rente créée par la collectivité, et le patrimoine improductif}

Une ligne de métro ou un lycée enrichissent les propriétaires des terrains alentour. Arnott et
Stiglitz (1979) donnent à l'intuition d'Henry George une forme précise : dans une ville de taille optimale, la somme des rentes foncières dues à la
localisation égale la dépense en biens publics locaux, si bien qu'un impôt unique sur ces rentes
suffit à la financer. Les auteurs doutent que ce théorème décrive les villes réelles ; il fonde
néanmoins l'idée qu'une LVT locale fait payer un équipement par ceux qu'il enrichit.

Le second argument porte sur le patrimoine. La hausse des patrimoines rapportés au revenu est pour
l'essentiel une hausse du prix du sol : Knoll, Schularick et Steger (2017) attribuent au terrain
environ 80~\% de la hausse des prix réels des logements depuis 1950 dans quatorze pays.  Taxer le sol plutôt que le
capital réoriente l'épargne vers le capital productif et élève les salaires (Stiglitz, 2015) ;
dans le modèle de Bonnet, Chapelle, Trannoy et Wasmer (2021), un impôt foncier sans impôt sur le
capital atteint l'optimum de premier rang. Kumhof, Tideman, Hudson et
Goodhart (2021) chiffrent à +15~\% le PIB américain d'un transfert massif de la fiscalité vers le
sol ; le chiffre sort d'un modèle calibré sans contrepartie empirique, et les auteurs le jugent
eux-mêmes « extrêmement élevé ».

\section{La France : beaucoup de terrain, une taxe foncière aveugle}\label{sec:france}

\subsection{Le retour de la terre}

Les terrains sous les bâtiments valaient l'équivalent de 50 à 75~\% du PIB des années 1970 à la fin
des années 1990. Ils en valent 229~\% fin 2024, après un pic à 296~\% en 2020, soit 6~729~Md€, dont
4~043~Md€ détenus par les ménages (figure~\ref{fig:foncier}). Les logements eux-mêmes, estimés au
coût de remplacement, sont passés de 147 à 204~\% du PIB : les trois quarts de la hausse de la
valeur de l'immobilier résidentiel rapportée au PIB depuis 1998 sont attribués au sol. La part du
terrain dans la valeur des biens des ménages est passée de 22~\% en 1998 à 56~\% en 2007, puis à
45~\% en 2024.

\figbloc{lvt_foncier_pib}{Valeur des terrains bâtis et des logements, et part du terrain dans
l'immobilier des ménages, 1978-2024}{fig:foncier}

\subsection{Que mesure la « valeur des terrains » ?}

Personne n'observe la valeur du sol sous un immeuble : on ne vend jamais l'un sans l'autre. L'Insee
la calcule par différence. La valeur totale du parc résidentiel est calée sur les enquêtes
Logement, dont la dernière utilisée date de 2013, puis prolongée ; on en retranche la valeur des
bâtiments, estimée en cumulant les investissements passés diminués d'une usure calculée sur une
durée de vie de 70 ans ; le reste est le terrain (Insee, fiche méthodologique de la base 2020).
Aucune marge d'incertitude n'est publiée, et la refonte complète du
partage entre bâti et terrain est reportée à la base de 2030.

Trois conséquences en découlent. D'abord, le terrain absorbe toutes les erreurs faites sur le
bâti : une valeur du bâti surestimée ou sous-estimée de 20~\% déplace la part du terrain d'environ dix points
(figure~\ref{fig:residu}), et le niveau de 2023 a été révisé de 9~\% (de 6~346 à 6~918~Md€) entre
deux publications de l'Insee. Le guide d'Eurostat et de l'OCDE (2015) signale même des valeurs
résiduelles négatives au Danemark en 1995-1996. Ensuite, le terrain absorbe toute la volatilité des prix : en 2023,
sa valeur a baissé de 8,9~\% quand celle du bâti augmentait de 3,0~\%.

Enfin, le résidu ne mesure pas l'assiette d'une LVT, qui est le terrain « comme s'il était
nu ». Le résidu attribue au sol tout ce qui, dans le prix d'un bien, dépasse le coût d'un bâtiment neuf
équivalent ; pour un immeuble vétuste promis à la démolition, il sous-estime la valeur du terrain
nu, et les deux mesures ne coïncident à peu près que pour les bâtiments de moins de vingt ans
(Davis, Larson, Oliner et Shui, 2021). Les comparaisons internationales souffrent du même défaut :
la France affiche des terrains de toute nature valant 293~\% du PIB en 2024 contre 174~\% en
Allemagne, mais l'Allemagne les mesure directement, à partir des valeurs de zone du sol, quand la
France procède par différence.

La hausse elle-même est toutefois solide, car on peut la lire sans passer par le résidu. Entre
1998 et 2025, le prix des logements anciens a été multiplié par 3,0, le coût de la construction
par 2,0, les loyers et les prix à la consommation par 1,5 (figure~\ref{fig:indices}). L'écart
entre le prix d'un logement et le coût de sa reconstruction, que les comptes attribuent au sol,
s'est donc bien creusé. Mais il ne traduit pas une hausse de ce que les logements rapportent : le
rapport entre prix et loyers a doublé, un mouvement qui coïncide avec la baisse des taux d'intérêt
jusqu'en 2021. La valeur du sol a bien plus augmenté que ce qu'il rapporte chaque année. Une taxe proportionnelle à la valeur suivrait ces mouvements de taux.

L'agrégat des comptes nationaux vaut donc pour l'ordre de grandeur et la tendance ; il ne
fournit ni un niveau précis, ni une assiette. Une LVT aurait besoin de valeurs par parcelle,
estimées directement sur les ventes.

\figbloc{lvt_prix_cout_loyers}{Prix des logements anciens, coût de la construction, loyers et prix
à la consommation, 1998-2025}{fig:indices}

\figbloc{lvt_residu}{Variation annuelle de la valeur des terrains et du bâti des ménages, et part
du terrain selon l'estimation du bâti}{fig:residu}

\subsection{Une valeur qui se joue à quelques kilomètres près}

En 2024, le prix médian d'un terrain à bâtir allait de 17~€ le mètre carré en Corrèze à 771~€ dans
le Val-de-Marne, un rapport de 45 (figure~\ref{fig:prix}). La part du terrain dans le prix d'une
maison varie dans les mêmes proportions : environ 10~\% en Corrèze, plus de 60~\% en Haute-Savoie,
dans le Rhône ou en Seine-Saint-Denis. Ces ordres de grandeur rejoignent ceux de Combes, Duranton
et Gobillon (2019), pour qui le terrain représente 36~\% du coût d'une maison neuve en France et
dont le prix croît trois fois plus vite que celui des logements avec la taille de la ville. À l'intérieur d'une même région, le prix du sol suit la
distance au centre : en Île-de-France, il passe d'environ 1~400~€/m² à moins de 10~km de
Notre-Dame à moins de 100~€/m² au-delà de 70~km (figure~\ref{fig:gradient}). C'est le mécanisme de
la rente de localisation : on paie le sol pour le temps de trajet qu'il épargne et les emplois,
écoles et services qu'il met à portée.

\figbloc{lvt_prix_terrains}{Prix des terrains à bâtir et part du terrain dans le prix des
maisons, par département, 2024}{fig:prix}

\figbloc{lvt_gradient_idf}{Prix des terrains à bâtir en Île-de-France selon la distance au centre
de Paris, 2024}{fig:gradient}

\subsection{Une taxe foncière qui taxe le confort de 1970}

La France taxe déjà beaucoup l'immobilier : les impôts périodiques sur la propriété immobilière
(taxes foncières, taxe d'habitation, cotisation foncière des entreprises, impôt sur la fortune
immobilière) y représentaient 1,9~\% du PIB en 2023, le premier niveau de l'Union européenne
(moyenne de l'OCDE : 1,0~\% ; Fipeco, 2025). La taxe foncière sur les propriétés bâties a rapporté 42~Md€ en 2024 ;
les droits de mutation, payés à chaque vente, 13,0~Md€ aux départements et 2,7~Md€ aux communes
en 2023. Le débat porte donc sur l'assiette, pas sur le niveau.

Cette assiette est la valeur locative cadastrale, le loyer que le bien aurait rapporté au
1\textsuperscript{er}~janvier~1970 : chaque logement a alors été classé dans l'une de huit
catégories, de « somptueux » à « délabré », puis évalué au mètre carré, majoré pour le confort. Ces
valeurs n'ont été actualisées qu'en 1980, puis revalorisées d'un coefficient uniforme. L'impôt porte sur la moitié de cette valeur, au taux
voté par la commune et l'intercommunalité (39,7~\% en moyenne en 2024). Un immeuble ancien du
centre de Paris, vétuste en 1970 et rénové depuis, reste classé comme il l'était ; un grand
ensemble de Seine-Saint-Denis, moderne en 1970, garde un classement élevé.

L'Insee a mesuré le résultat sur les données fiscales de 2017 (André et Meslin, 2023), en
rapportant la taxe de chaque ménage à une valeur de marché estimée de ses biens : la taxe
représente 0,34~\% de cette valeur en moyenne, plus de 0,5~\% pour les plus petits patrimoines,
0,22~\% pour le centile le plus doté ; environ 0,1~\% de la valeur des logements à Paris, souvent
plus de 0,4~\% en Seine-Saint-Denis. Sur 4~661 communes, la taxe pèse
0,54~\% de la valeur des logements dans les 5~\% les moins chères et 0,20~\% dans les 5~\% les plus
chères (figure~\ref{fig:tf}). Rapportée au revenu, elle reste progressive : les 10~\% de ménages les
plus aisés en paient 27~\%. La révision des valeurs d'habitation, votée en 2019 pour 2026, a été
reportée à 2028 puis à 2031 ; elle se ferait sur les loyers, pas sur la valeur vénale ni sur le
terrain. Elle augmenterait de plus de 9~\% les valeurs des ménages les plus aisés (Chapelle, Fabre
et Lallemand, 2020) : c'est ce transfert qui a fait abandonner la révision de 1990.

\figbloc{lvt_tf_regressive}{Taxe foncière rapportée à la valeur des logements, selon le prix au
mètre carré de la commune}{fig:tf}

\subsection{Ordres de grandeur}

Les terrains bâtis détenus hors administrations publiques valaient 5~826~Md€ fin 2024. Une taxe
qui rapporterait les 42~Md€ de la taxe foncière bâtie aurait un taux de 0,72~\% de la valeur du
terrain ; en remplaçant aussi les droits de mutation de 2023, le taux serait proche de 1,0~\%. La
charge se déplacerait fortement entre biens. La maison médiane vendue en 2024 en Corrèze
(118~000~€, terrain d'environ 10~\% de la valeur) et celle de Haute-Savoie (461~000~€, terrain
d'environ 65~\%) paient aujourd'hui à peu près la même taxe, 1~040 et 930~€. Avec une taxe de
0,57~\% du terrain, le taux qui aurait rapporté en 2017 autant que la taxe foncière des ménages sur
leurs logements, la première paierait environ 70~€ et la seconde 1~700~€. Les montants sont
approximatifs ; le passage d'un rapport de 1 à 1 à un rapport de 1 à 25 ne l'est pas. Un
appartement parisien de 500~000~€, dont le terrain représenterait 75~\% de la valeur (hypothèse
qu'aucune donnée publique ne permet de vérifier), paie aujourd'hui environ 730~€, 0,15~\% de sa
valeur après la hausse de taux de 2023 ; il paierait 2~130~€.

\section{Ce que mesurent les études, transposé à la France}\label{sec:etudes}

Pour chaque étude, on applique l'élasticité estimée à une grandeur française par une règle de
trois. Ces transpositions supposent que l'effet mesuré ailleurs vaudrait en France et resterait
proportionnel au-delà des variations observées : ce sont des ordres de grandeur, pas des
prévisions.

\subsection{Les prix et les loyers : qui paie}

\textbf{Capitalisation.} Les impôts fonciers ordinaires sont en grande partie capitalisés. En
Angleterre, de part et d'autre des frontières entre collectivités, les écarts de \emph{council tax}
se retrouvent dans les prix comme si les acheteurs actualisaient les impôts futurs à 3,7~\%, soit une
capitalisation presque complète (Koster et Pinchbeck, 2022). Les travaux plus anciens trouvaient deux tiers (Oates, 1969) ou 62~\%
(Palmon et Smith, 1998). En France, la suppression de la taxe d'habitation a été capitalisée
à 17 à 48~\% dans les prix selon le taux d'actualisation retenu, et répercutée pour moitié environ
dans les loyers (Bach, Bozio et
al., IPP, 2023).

\textbf{Le cas danois.} La seule expérience naturelle sur une taxe purement foncière est la fusion
des communes danoises en 2007 : l'harmonisation des taux a fait monter ou baisser durablement la
taxe de propriétaires voisins, sans rapport avec leur marché local. Høj, Jørgensen et Schou (2018), sur 211~000 ventes de maisons, trouvent qu'un point pour
mille de taxe (rapportée à la valeur du bien) réduit le prix de 4,3~\% : capitalisation complète,
au taux implicite de 2,3~\%, et déjà visible dès l'annonce de la réforme en 2004. Nielsson,
Wroblewski et Yding (2024) reprennent la réforme avec les ventes répétées des mêmes logements et
un second instrument, l'intensité foncière historique de chaque parcelle. Ils trouvent un effet
nul et précis, qui exclut une capitalisation complète pour tout taux d'actualisation inférieur à
8~\%. Ils montrent que le résultat de Høj et al. disparaît quand on tient compte des résidences de
vacances, plus nombreuses là où la taxe a monté et dont les prix ont divergé au même moment.
Ils ne trouvent pas d'effet sur la construction ni sur la mobilité, mais les propriétaires
âgés quittent les communes où la taxe a augmenté. Si les prix ne
baissent pas, la taxe est payée par chaque propriétaire successif, et par les locataires si elle
est répercutée : peu distorsive, elle devient plus régressive. Les auteurs l'expliquent par le contrôle des loyers
danois, qui autorise le bailleur à répercuter les hausses d'impôt ; l'explication ne rend pas
compte de l'absence d'effet sur les maisons occupées par leur propriétaire.

\textbf{Répercussion sur les loyers.} En Allemagne, où l'impôt foncier fait partie des charges
refacturables au locataire, Löffler et Siegloch (2024) mesurent sur 5~500 hausses de taux
communaux qu'environ 83~\% de la hausse se retrouve dans les loyers après trois ans. À Berkeley,
la répercussion est de 50 à 89~\% (Baker, 2025). En France, la taxe foncière n'est pas récupérable sur le locataire (le
décret de 1987 sur les charges récupérables ne la mentionne pas), mais le bailleur peut relever le
loyer à la relocation.

La figure~\ref{fig:reglesprix} applique ces élasticités aux biens types de la
section~\ref{sec:france}, dans une bascule à recette constante. Avec l'élasticité de Høj et al.,
l'appartement parisien perdrait 12~\% de sa valeur et la maison haut-savoyarde 7~\% ; avec celle de
Koster et Pinchbeck, 8 et 4~\% ; avec Nielsson et al., au plus 3,5 et 2~\%, pour une estimation
centrale nulle. Pour la maison corrézienne, dont la taxe baisserait de 8~‰ de sa valeur, les
transpositions sortent du domaine observé. Si les biens sont loués et que la répercussion
allemande s'applique, le loyer annuel de l'appartement parisien augmenterait d'environ 1~160~€ et
celui de la maison corrézienne baisserait de 810~€.

\figbloc{lvt_regles_prix}{Effet d'une bascule vers une taxe sur le terrain sur le prix et le loyer
de biens types, selon les élasticités publiées}{fig:reglesprix}

\subsection{La construction : des effets positifs et modestes}

Les preuves viennent de Pennsylvanie. À Pittsburgh, qui a porté en 1979-1980 son taux municipal sur
le terrain à plus de cinq fois celui sur le bâti, la valeur des permis de construire a augmenté de
70~\% entre 1960-1979 et 1980-1989, quand elle baissait de 14~\% dans quinze villes comparables ; les
auteurs attribuent l'essentiel du boom à une pénurie de bureaux et donnent à la taxe un « rôle
d'appoint » (Oates et Schwab, 1997). Sur l'ensemble des villes concernées, Plassmann et Tideman (2000) trouvent de
l'ordre de 3 à 4~\% de construction résidentielle en plus (selon la synthèse de Dye et England,
2010) ; Banzhaf et Lavery (2010), une croissance du nombre de logements supérieure de 0,6 point par
décennie pour un taux foncier double du taux sur le bâti, de 2 points pour un taux quadruple, par
davantage de logements plutôt que des logements plus grands. Dye et England jugent ces
résultats « le plus souvent peu concluants ». L'effet de l'impôt sur le bâti est mieux établi : dans le New Hampshire, un allègement de l'impôt foncier égal à
15~\% des recettes locales a accru la construction de 11 à 22~\%, là où l'offre est souple ; près
de Boston, où elle ne l'est pas, il s'est retrouvé dans les prix (Lutz, 2015).

Transposés à la France (figure~\ref{fig:reglesconstruction}), ces effets donnent 8~000 à
11~000 logements commencés de plus par an pour Plassmann et Tideman, sur 263~000 en 2025, et
19~000 à 63~000 pour Banzhaf et Lavery, appliqués à un parc de 31,5~millions de résidences
principales. Le second majore l'effet national, car il inclut des logements qui se seraient construits
dans d'autres quartiers. La LVT serait un appoint à la politique du logement, pas un choc d'offre.

\figbloc{lvt_regles_construction}{Logements supplémentaires par an si les effets mesurés en
Pennsylvanie se reproduisaient en France}{fig:reglesconstruction}

\subsection{Les droits de mutation : là où les dégâts sont les mieux mesurés}

Les preuves les plus solides portent sur ce qu'une LVT pourrait remplacer. Un droit de mutation
taxe le fait de changer de logement. En Allemagne, un point de droit supplémentaire réduit les
transactions d'environ 7~\% (Fritzsche et Vandrei, 2019) et les prix de 3~\% (Dolls et al., 2025) ;
à Toronto, un droit de 1,1~\% a réduit les ventes de 15~\%, avec un dollar de perte sèche pour huit
de recette (Dachis, Duranton et Turner, 2012) ; au Royaume-Uni, deux points de \emph{stamp duty}
réduisent les déménagements de 37~\% (Hilber et Lyytikäinen, 2017). La seule estimation française
nuance ces résultats : la hausse de 0,7 point de 2014 a avancé les ventes (+26~\% le mois précédent)
puis les a freinées trois mois, sans effet au-delà (Bérard et Trannoy, 2018).

Appliquées aux 778~000 ventes de logements recensées dans DVF en 2024, les élasticités étrangères
donnent 54~000 à 106~000 ventes de plus par an pour un point de droit en moins
(figure~\ref{fig:reglesdmto}). Prolongées jusqu'à la suppression des droits (5,8 points sur
l'ancien), elles donneraient +41 à +79~\% de ventes, très au-delà de ce qui a été observé ; le
modèle estimé sur données françaises de Bontemps, Cherbonnier et Magnac (2023) trouve +21~\% pour
un passage de 14 à 6~\% de coûts de transaction compensé par un impôt sur la propriété, avec un
gain de bien-être de 0,2~\% seulement, car les contraintes de crédit, plus que les droits, bornent
les choix des ménages.

\figbloc{lvt_regles_dmto}{Effet d'une baisse des droits de mutation sur les ventes de logements,
selon les élasticités publiées}{fig:reglesdmto}

\section{Ce que donne une mise en place progressive}\label{sec:transition}

\subsection{Le mécanisme : étaler la taxe n'étale pas la perte}

On pourrait croire qu'une taxe montée en charge sur quinze ans est quinze fois moins douloureuse.
Ce n'est pas le cas si les acheteurs anticipent : le jour où la trajectoire est votée, tous les
impôts futurs, y compris ceux de la quinzième année et des suivantes, entrent dans le prix du
terrain. La montée en charge ne retire de la perte que les impôts non perçus pendant la
transition. Pour une taxe de 1~\% de la valeur du terrain et un taux d'actualisation de 3~\%, le
prix baisse de 25~\% à l'annonce si la taxe s'applique tout de suite, de 22~\% si elle monte en dix
ans, de 19~\% en vingt ans (figure~\ref{fig:trajectoire}). La part de la perte supportée à
l'annonce est de 88~\% pour une transition de dix ans, 77~\% pour vingt ans, encore 67~\% pour
trente ans. Le prix glisse ensuite jusqu'à $-25$~\%, mais ce glissement ne lèse personne, car celui
qui achète pendant la transition paie moins cher.

\figbloc{lvt_trajectoire}{Prix d'un terrain après l'annonce d'une taxe de 1~\% de sa valeur,
selon la durée de montée en charge}{fig:trajectoire}

Ce que les propriétaires ne perdent pas à l'annonce est exactement la valeur actualisée des
recettes auxquelles l'État renonce pendant la transition : la montée en charge est une subvention
aux propriétaires du moment, payée par les autres contribuables. Ce qu'elle protège vraiment est la
trésorerie : une taxe qui doit passer de 900 à 1~700~€ par an n'augmente que de 50~€ par an sur
quinze ans.

\subsection{Ce que cela donne pour des biens types}

Dans une bascule à recette constante, il y a autant de gains que de pertes en capital, et la
transition les réduit dans la même proportion (figure~\ref{fig:bienstypes}). La maison corrézienne,
dont la facture passerait de 1~040 à 70~€, gagnerait 28~\% de valeur si la bascule était immédiate,
23~\% si elle était étalée sur quinze ans. La maison rhodanienne ne bouge pas : son terrain pèse
beaucoup, mais sa taxe actuelle aussi. La maison haut-savoyarde perdrait 5,5~\% (4,5~\% sur quinze
ans), l'appartement parisien 9~\% (8~\%), soit environ 47~000~€. Une telle perte est de l'ordre d'une
année de variation des prix ; elle n'est pas confiscatoire, mais elle se concentre sur les
acquéreurs récents, dont le prêt couvre en moyenne 85~\% de la valeur du bien pour un premier achat (Comité
consultatif du secteur financier, 2026) : avec un apport de 15~\%, une baisse de 9~\% du prix efface
plus de la moitié des fonds propres, alors que la dette ne bouge pas.

\figbloc{lvt_biens_types}{Variation du prix de biens types à l'annonce d'une bascule vers une
taxe sur le terrain, selon la durée de montée en charge}{fig:bienstypes}

\subsection{Ce que montrent les transitions réelles}

Le calcul suppose que les acheteurs anticipent ; ils le font mal. Au
Michigan, l'impôt foncier d'un bien est plafonné tant que son propriétaire ne change pas, puis
réaligné après la vente : les acheteurs ne voient pas la hausse programmée et surpaient leur
logement d'environ 10~000 dollars en moyenne (Bradley, 2017). En Suède, la réforme de 2008 n'a
eu d'effet sur les prix ni à l'annonce ni à l'application, sauf pour les biens les plus chers
(Elinder et Persson, 2017) ; au Danemark, à l'inverse, Høj et al. voient les prix bouger dès
l'annonce. Si l'anticipation est incomplète, la montée en charge protège les propriétaires du
moment aux dépens de ceux qui achètent pendant la transition. D'où une règle simple : afficher, à
chaque vente, la taxe que le bien paiera à terme.

Les expériences étrangères montrent aussi ce qui fait échouer une transition. Pittsburgh n'a pas
renoncé à son double taux par rejet du principe, mais après une réévaluation brutale : le comté
avait gelé les valeurs de 1996 à 2001, et le rattrapage imposé par un juge a fait bondir de 81~\%
la valeur foncière imposable en une année (Hughes, 2006). En Nouvelle-Zélande, la part des communes qui
taxent la valeur du sol est tombée de 47~\% en 2007 à 29~\% en 2019, au profit de la valeur
bâtie.

Trois dispositifs amortissent le choc. Le plafonnement annuel
attaché au bien, comme au Danemark depuis 2024 (la taxe sur le terrain d'un bien ne peut croître
de plus de 4,75~\% par an) ou en Norvège (au plus un point pour mille de plus par an), lisse les
factures sans dissuader de vendre. L'avantage attaché au propriétaire, comme le rabais danois qui
garantit aux propriétaires de 2023 de ne pas payer davantage mais disparaît à la vente, crée au
contraire un verrou : en Californie, où l'avantage fiscal de la Proposition~13 devient
transférable à 55 ans, la part des ménages qui déménagent en cinq ans passe de 4 à 5,2~\% à cet âge
(Ferreira, 2010). Le report de paiement, réglé
à la vente ou à la succession, protège les ménages modestes sans éroder l'assiette, mais il n'est
utilisé que s'il est automatique : aux États-Unis, au milieu des années 1990, moins de 1~\% des ménages
éligibles y recouraient (Shan, 2010). Pour une bascule française, la combinaison la plus défendable serait une montée en
charge inscrite dans la loi, un plafond annuel par bien et un report automatique pour les
propriétaires à faibles revenus.


\section{Les autres difficultés}

\subsection{Estimer la valeur de chaque terrain}

Il faut une valeur par parcelle, contestable par le contribuable. Trois méthodes existent : les ventes de terrains nus, rares là où le sol est cher (27~000 ventes de
terrains à bâtir dans DVF en 2024, presque aucune à Paris) ; le résidu, qui reporte sur le sol
toutes les erreurs faites sur le bâti ; des valeurs de zone estimées sur l'ensemble des ventes et
appliquées à chaque parcelle, la solution allemande et danoise. La France a la matière première
(DVF, cadastre, loyers déclarés depuis 2023) mais pas l'institution : personne n'y estime de
valeurs foncières de zone, et la DGFiP objecte qu'un bien ne se vend en moyenne que tous les
dix-neuf ans (Conseil des prélèvements obligatoires, 2023). Une valeur imparfaite suffit à l'efficacité, puisque le propriétaire ne
peut pas modifier son impôt par ses décisions (\emph{Mirrlees Review}, 2011) ; elle ne suffit pas à
l'acceptation, comme l'a montré Pittsburgh.

\subsection{Les revenus, les communes, le zonage}

Le transfert vers les zones chères est progressif au regard du patrimoine (les 10~\% de ménages les
mieux dotés détiennent environ 44~\% du patrimoine immobilier) mais pas toujours au regard du
revenu : un retraité modeste propriétaire d'une maison achetée il y a quarante ans dans une zone
devenue chère verrait sa taxe multipliée sans que son revenu change. Le report de paiement, s'il est
automatique, règle ce problème (section~\ref{sec:transition}).

La taxe foncière est la principale ressource fiscale des communes depuis la suppression de la
taxe d'habitation sur les résidences principales, et ce sont elles qui en votent le taux. Changer
l'assiette redistribue les bases entre communes autant qu'entre contribuables. La révision des
valeurs des locaux professionnels, appliquée en 2017, a montré la méthode : coefficient de
neutralisation qui gèle la part de chaque commune, lissage sur dix ans, plafonnement des écarts.

Enfin, la valeur taxée dépend du plan local d'urbanisme : rendre un terrain constructible
multiplie sa valeur et sa taxe. La LVT ne crée pas d'offre là où les règles interdisent de
construire, et elle donne aux propriétaires une raison de plus de s'opposer à l'extension des
droits à construire, qui alourdirait leur impôt avant d'enrichir leur patrimoine.

\clearpage
\section{Options pour la France}

Les taux sont calculés à recette constante, sur la valeur des terrains fin 2024.

\begin{center}\small
\begin{tabularx}{\linewidth}{@{}>{\raggedright\arraybackslash}p{3.1cm}>{\raggedright\arraybackslash}X>{\raggedright\arraybackslash}p{4.3cm}@{}}
\toprule
Option & Ce qui change & Ce qu'il faut construire \\
\midrule
1. Valeur de marché & La taxe foncière porte sur la valeur vénale estimée du bien entier
(environ 0,3~\% en moyenne). Corrige la régressivité ; ne change pas l'incitation à construire. &
Estimation statistique de chaque bien (DVF, cadastre) ; actualisation régulière ; neutralisation
communale. \\[3pt]
2. Double taux & Une part croissante de la base est le terrain seul (taux deux à trois fois plus
élevé sur le sol que sur le bâti). Allège l'impôt sur la construction. & Valeurs de zone du
terrain, publiées et contestables, sur le modèle allemand. \\[3pt]
3. LVT pure & La taxe foncière bâtie est remplacée par une taxe sur le terrain : environ 0,7~\% de
sa valeur par an. & Idem, plus un report de paiement pour les propriétaires à faibles revenus et
une péréquation entre communes. \\[3pt]
4. LVT et droits de mutation & La LVT remplace aussi les droits de mutation : environ 1~\% de la
valeur du terrain. & Idem, plus une compensation des départements. \\
\bottomrule
\end{tabularx}
\end{center}

L'option~1 est un préalable et la plus défendable à court terme : elle corrige le défaut le mieux
documenté et produit l'infrastructure (une valeur de marché par bien) dont une LVT aurait besoin.
L'option~2 introduit la LVT à petite dose ; le Bade-Wurtemberg montre qu'elle tient juridiquement. Les options~3 et~4 produisent le gain
d'efficacité le plus grand, mais aussi les transferts décrits en sections~\ref{sec:france}
et~\ref{sec:transition} dans toute leur ampleur, sur une assiette que la France ne sait pas encore
mesurer. Trois questions restent à trancher : viser la recette constante ou financer une baisse
d'autres impôts ; réforme nationale à taux unique ou laissée aux communes ; et, tant que le résultat
danois n'est pas tranché, accepter de ne pas savoir qui paierait vraiment.

\appendix
\section*{Annexe : sources et méthode des calculs}
\addcontentsline{toc}{section}{Annexe}

\small
\textbf{Valeur des terrains.} Insee, comptes de patrimoine annuels en base 2020, tableaux 8.201
(économie nationale) et 8.210 (ménages), ligne « terrains supportant des bâtiments et des ouvrages
de génie civil » ; PIB à prix courants (Insee, base 2020). Fin 2024 : 6~729~Md€ pour l'économie
nationale, dont 903~Md€ pour les administrations publiques (Insee Première n°~2081). Les taux
neutres rapportent la taxe foncière bâtie 2024 (42~Md€, DGCL, note sur le fichier REI) et les
droits de mutation 2023 des départements et des communes (13,0 et 2,7~Md€, réponse ministérielle à
la question écrite n°~4482, 2025) aux 5~826~Md€ de terrains hors administrations. Le taux de 0,57~\%
rapporte la taxe foncière des ménages sur leurs logements en 2017 (19,2~Md€, André et Meslin,
2023) à la valeur des terrains bâtis des ménages fin 2017 (3~377~Md€).

\textbf{Prix des terrains et des maisons.} DGFiP, Demandes de valeurs foncières 2024, version
géolocalisée d'Etalab. Terrains à bâtir : ventes dont toutes les parcelles sont classées « terrain
à bâtir » au cadastre, sans local, de 100 à 5~000~m² (27~353 ventes). Maisons : ventes d'une seule
maison, sans appartement ni local d'activité, de 30 à 400~m² habitables et 50 à 10~000~m² de
terrain (304~799 ventes). Part du terrain : prix médian du m² de terrain à bâtir du département
multiplié par la surface médiane de parcelle des maisons vendues, rapporté au prix médian des
maisons. Ce calcul sous-estime probablement la part du terrain dans les villes, où les maisons
sont sur des terrains plus chers que les terrains à bâtir vendus en périphérie, et ne vaut que
comme ordre de grandeur. L'Alsace-Moselle est absente de DVF.

\textbf{Taxe foncière et prix.} Taxe foncière sur les logements rapportée à leur valeur de marché
estimée, par commune, en 2017 : données de la figure~5a de l'Insee Analyses n°~91. Prix médian au
m² des ventes d'un seul logement en 2024 (DVF), pour les 4~661 communes comptant au moins 20 ventes.
Élasticité du ratio taxe/valeur au prix au m², estimée en coupe sur ces communes : $-0{,}54$. Taux
départementaux de l'exemple : médiane des communes du département pondérée par le nombre de ventes.

\textbf{Indices et volumes.} Insee : indice Notaires-Insee des prix des logements anciens (France
métropolitaine, brut), indice du coût de la construction, indice des prix à la consommation base
2015 (loyers effectivement payés ; ensemble hors tabac), en moyennes annuelles ; logements commencés
(Sitadel, cumul sur douze mois à fin décembre) ; résidences principales (estimations annuelles du
parc). Ventes de logements : mutations « vente » de DVF 2024 comportant au moins une maison ou un
appartement (777~629). Gradient francilien : 1~159 ventes de terrains à bâtir à moins de 80~km de
Notre-Dame, distance à vol d'oiseau.

\textbf{Mise en place progressive.} Un terrain rapporte un loyer net constant $R$ ; la taxe de
l'année $k$ vaut $t_k P_k$, avec $t_k$ croissant linéairement de 0 au taux final en $N$ années. Le
prix vérifie $P_{k-1} = \left(R + (1-t_k) P_k\right)/(1+r)$, résolu à rebours à partir du régime
permanent $P = R/(r+t)$, avec $r = 3$~\%. Pour les biens types, la variation de la facture
annuelle $\Delta T$ est supposée fixe ; la variation de prix à l'annonce vaut alors
$-\varphi_N \Delta T / r$, où $\varphi_N$ est la part de la valeur actualisée d'une hausse
permanente qui reste due avec une montée linéaire sur $N$ ans (0,88 pour dix ans, 0,82 pour quinze).

\textbf{Règles de trois.} Prix : variation de la taxe rapportée à la valeur du bien, en points pour
mille, multipliée par la baisse de prix par point pour mille estimée (Høj et al. : 4,26~\% ;
Koster et Pinchbeck : 1/3,7~\%, soit 2,7~\% ; Nielsson et al. : au plus 1,25~\%). Loyers : 83~\%
de la variation de taxe (Löffler et Siegloch), 50 à 89~\% (Baker). Ventes : 7~\% par point de
droit (Fritzsche et Vandrei), 15~\% pour 1,1 point (Dachis et al.), 37~\% de mobilité pour 2 points
(Hilber et Lyytikäinen). Construction : 3 à 4~\% des logements commencés (Plassmann et Tideman) ;
0,6 à 2 points de croissance décennale du parc (Banzhaf et Lavery, rapport des taux de 2 et 4).

Les scripts \texttt{extract\_lvt.py}, \texttt{extract\_complements.py} et
\texttt{graphiques\_lvt.py} du dossier \texttt{fiscalite/lvt} refont tous les calculs à partir
des fichiers publics.

\section*{Références citées}
\begin{list}{}{\leftmargin=1.2em\itemindent=-1.2em\itemsep=1pt\parsep=0pt}
\item André, M. et O. Meslin (2023), « Les ménages au patrimoine immobilier le plus élevé paient
relativement moins de taxe foncière que les ménages les moins dotés », \emph{Insee Analyses} n°~91.
\item Arnott, R. et J. Stiglitz (1979), « Aggregate Land Rents, Expenditure on Public Goods, and
Optimal City Size », \emph{Quarterly Journal of Economics} 93(4), 471-500.
\item Bach, L., A. Bozio, P. Dutronc-Postel, É. Fize, A. Guillouzouic et C. Malgouyres (2023),
\emph{Évaluation de la réforme de la taxe d'habitation}, rapport IPP n°~48.
\item Baker, S. (2025), « Property Tax Pass-Through to Renters: A Quasi-Experimental Approach »,
Federal Reserve Bank of Philadelphia WP 25-41.
\item Banzhaf, H.~S. et N. Lavery (2010), « Can the Land Tax Help Curb Urban Sprawl? Evidence from
Growth Patterns in Pennsylvania », \emph{Journal of Urban Economics} 67(2), 169-179.
\item Bentick, B. (1979), « The Impact of Taxation and Valuation Practices on the Timing and
Efficiency of Land Use », \emph{Journal of Political Economy} 87(4), 859-868.
\item Bérard, G. et A. Trannoy (2018), « The impact of the 2014 increase in the real estate transfer
taxes on the French housing market », \emph{Économie et Statistique} 500-501-502, 179-200.
\item Bonnet, O., G. Chapelle, A. Trannoy et É. Wasmer (2021), « Land is back, it should be taxed,
it can be taxed », \emph{European Economic Review} 134, 103696.
\item Bontemps, C., F. Cherbonnier et T. Magnac (2023), « Reducing transaction taxes on housing in
highly regulated economies », TSE Working Paper 23-1486.
\item Bradley, S. (2017), « Inattention to Deferred Increases in Tax Bases: How Michigan Home Buyers
Are Paying for Assessment Limits », \emph{Review of Economics and Statistics} 99(1), 53-66.
\item Brueckner, J. et H.-A. Kim (2003), « Urban Sprawl and the Property Tax », \emph{International
Tax and Public Finance} 10(1), 5-23.
\item Calvo, G., L. Kotlikoff et C. Rodriguez (1979), « The Incidence of a Tax on Pure Rent: A New
(?) Reason for an Old Answer », \emph{Journal of Political Economy} 87(4), 869-874.
\item Chapelle, G., B. Fabre et C. Lallemand (2020), « Révision des valeurs locatives sur les locaux
d'habitation », rapport IPP n°~29.
\item Combes, P.-P., G. Duranton et L. Gobillon (2019), « The Costs of Agglomeration: House and Land
Prices in French Cities », \emph{Review of Economic Studies} 86(4), 1556-1589.
\item Comité consultatif du secteur financier (2026), \emph{Le modèle français de financement de
l'acquisition de la résidence principale}.
\item Conseil des prélèvements obligatoires (2023), \emph{Pour une fiscalité du logement plus
cohérente}.
\item Dachis, B., G. Duranton et M. Turner (2012), « The effects of land transfer taxes on real
estate markets », \emph{Journal of Economic Geography} 12(2), 327-354.
\item Davis, M., W. Larson, S. Oliner et J. Shui (2021), « The price of residential land for
counties, ZIP codes, and census tracts in the United States », \emph{Journal of Monetary
Economics} 118, 413-431.
\item Dolls, M., C. Fuest, C. Krolage et F. Neumeier (2025), « Who bears the burden of real estate
transfer taxes? », \emph{Journal of Urban Economics} 145, 103717.
\item Dye, R. et R. England (2010), \emph{Assessing the Theory and Practice of Land Value Taxation},
Lincoln Institute of Land Policy.
\item Elinder, M. et L. Persson (2017), « House price responses to a national property tax
reform », \emph{Journal of Economic Behavior and Organization} 144, 18-39.
\item Eurostat et OCDE (2015), \emph{Eurostat-OECD compilation guide on land estimation}.
\item Feldstein, M. (1977), « The Surprising Incidence of a Tax on Pure Rent: A New Answer to an Old
Question », \emph{Journal of Political Economy} 85(2), 349-360.
\item Ferreira, F. (2010), « You can take it with you: Proposition 13 tax benefits, residential
mobility, and willingness to pay for housing amenities », \emph{Journal of Public Economics} 94(9-10),
661-673.
\item Fipeco (2025), « Les taxes foncières payées par les ménages », fiche du 29 septembre 2025.
\item Fritzsche, C. et L. Vandrei (2019), « The German real estate transfer tax: Evidence for
single-family home transactions », \emph{Regional Science and Urban Economics} 74, 131-143.
\item Hilber, C. et T. Lyytikäinen (2017), « Transfer taxes and household mobility »,
\emph{Journal of Urban Economics} 101, 57-73.
\item Høj, A., M. Jørgensen et P. Schou (2018), « Land Tax Changes and Full Capitalisation »,
\emph{Fiscal Studies} 39(2), 365-380.
\item Hughes, M.~A. (2006), \emph{Why So Little Georgism in America: Using the Pennsylvania Case
Files to Understand the Slow, Uneven Progress of Land Value Taxation}, Lincoln Institute of Land
Policy, WP06ZK1.
\item Insee (2026), « La refonte du compte de patrimoine non financier en base 2020 », fiche
méthodologique n°~24.
\item Knoll, K., M. Schularick et T. Steger (2017), « No Price Like Home: Global House Prices,
1870-2012 », \emph{American Economic Review} 107(2), 331-353.
\item Koster, H. et E. Pinchbeck (2022), « How Do Households Value the Future? Evidence from
Property Taxes », \emph{American Economic Journal: Economic Policy} 14(1), 207-239.
\item Kumhof, M., N. Tideman, M. Hudson et C. Goodhart (2021), « Post-Corona Balanced-Budget
Super-Stimulus: The Case for Shifting Taxes onto Land », CEPR DP 16652.
\item Löffler, M. et S. Siegloch (2024), « Welfare Effects of Property Taxation », ECONtribute
DP 331.
\item Lutz, B. (2015), « Quasi-experimental evidence on the connection between property taxes and
residential capital investment », \emph{American Economic Journal: Economic Policy} 7(1), 300-330.
\item Mills, D. (1981), « The Non-Neutrality of Land Value Taxation », \emph{National Tax Journal}
34(1), 125-129.
\item Mirrlees, J. et al. (2011), \emph{Tax by Design}, ch.~16, Oxford University Press.
\item Nielsson, U., C. Wroblewski et A. Yding (2024), « The Incidence and Efficiency of Land Value
Taxation », document de travail.
\item Oates, W. (1969), « The Effects of Property Taxes and Local Public Spending on Property
Values », \emph{Journal of Political Economy} 77(6), 957-971.
\item Oates, W. et R. Schwab (1997), « The Impact of Urban Land Taxation: The Pittsburgh
Experience », \emph{National Tax Journal} 50(1), 1-21.
\item OCDE (2022), \emph{Housing Taxation in OECD Countries}, OECD Tax Policy Studies n°~29.
\item Palmon, O. et B. Smith (1998), « New Evidence on Property Tax Capitalization », \emph{Journal
of Political Economy} 106(5), 1099-1111.
\item Plassmann, F. et T.~N. Tideman (2000), « A Markov Chain Monte Carlo Analysis of the Effect of
Two-Rate Property Taxes on Construction », \emph{Journal of Urban Economics} 47(2), 216-247.
\item Segú, M. (2020), « The impact of taxing vacancy on housing markets: Evidence from France »,
\emph{Journal of Public Economics} 185, 104079.
\item Shan, H. (2010), « Property taxes and elderly mobility », \emph{Journal of Urban Economics}
67(2), 194-205.
\item Stiglitz, J. (2015), « New Theoretical Perspectives on the Distribution of Income and Wealth
among Individuals: Part IV. Land and Credit », NBER WP 21192.
\item Tideman, T.~N. (1982), « A Tax on Land Value Is Neutral », \emph{National Tax Journal} 35(1),
109-111.
\item Wildasin, D. (1982), « More on the Neutrality of Land Taxation », \emph{National Tax Journal}
35(1), 105-108.
\end{list}
\end{document}
